\documentclass[11pt]{article}
\usepackage[paperwidth=160mm,paperheight=90mm,left=10mm,right=10mm,top=7mm,bottom=6mm]{geometry}
\usepackage[spanish]{babel}
\usepackage[T1]{fontenc}
\usepackage{xcolor}
\usepackage{hyperref}

\definecolor{Pino}{HTML}{143E37}
\definecolor{Mineral}{HTML}{007564}
\definecolor{Papel}{HTML}{F5F6F1}
\definecolor{Tinta}{HTML}{1A2E2A}
\pagecolor{Papel}
\color{Tinta}
\renewcommand{\familydefault}{\sfdefault}
\setlength{\parindent}{0pt}
\hypersetup{colorlinks, urlcolor=Mineral, pdftitle={Jaleco Sports — Agente de postventa}, pdfauthor={Grupo 5 — UNAL Manizales}}

\newenvironment{slide}[1]{%
  \thispagestyle{empty}%
  \begin{minipage}[t][0.97\textheight][t]{\textwidth}%
  \begingroup\setlength{\fboxsep}{6pt}%
  \colorbox{Pino}{\parbox{\dimexpr\linewidth-2\fboxsep\relax}{\color{white}\Large\bfseries #1}}%
  \endgroup\par\medskip
}{%
  \par\vfill\hrule\smallskip
  {\scriptsize Grupo 5 · Jaleco Sports\hfill\thepage\,/\,10}\par
  \end{minipage}\newpage
}

\newenvironment{block}[1]{%
  \par\medskip\textcolor{Mineral}{\bfseries #1}\par\smallskip
}{\par\medskip}

\begin{document}

\begin{slide}{Reclamos de postventa en Jaleco Sports}
  {\LARGE\bfseries Del mensaje libre al caso radicado}\par\medskip
  Jaleco Sports es una tienda ficticia de artículos deportivos. Los clientes describen un problema con su pedido y soporte debe identificar el motivo, reunir datos y aplicar la regla correspondiente.
  \begin{block}{Quiénes usan el sistema}
    El cliente completa su solicitud por chat. Soporte recibe un caso con pedido, motivo, evidencia pendiente y referencia de la regla.
  \end{block}
  \medskip Procesamiento de Lenguaje Natural · UNAL Manizales
\end{slide}

\begin{slide}{Proceso actual y costo por medir}
  \begin{enumerate}
    \item Interpretar el mensaje y pedir información faltante.
    \item Consultar el pedido y localizar la política vigente.
    \item Transcribir la solicitud y derivarla a soporte.
  \end{enumerate}
  \begin{block}{Medición prevista}
    Comparar el tiempo hasta un registro utilizable y los campos completos en el recorrido manual y asistido, sobre los mismos casos sintéticos.
  \end{block}
\end{slide}

\begin{slide}{Propuesta y alcance}
  \begin{itemize}
    \item El modelo distingue cambio por talla, defecto, pedido incompleto y retracto. Pide aclaración cuando hay ambigüedad.
    \item El workflow consulta el pedido y aplica una regla versionada con código.
    \item El cliente revisa la solicitud, confirma la radicación y descarga su constancia.
  \end{itemize}
  \begin{block}{Límite de la decisión}
    El resultado es preliminar. Los casos complejos pasan a soporte. Reembolsos, transporte y verificación visual quedan para trabajo futuro.
  \end{block}
\end{slide}

\begin{slide}{Recorrido de la demostración}
  \begin{enumerate}
    \item Identificar un pedido en la sesión A y recuperarlo en el siguiente turno.
    \item Repetir el seguimiento en la sesión B y comprobar que falta contexto.
    \item Mostrar datos del pedido y el cálculo de un plazo con un festivo.
    \item Intentar radicar sin confirmar, comprobar el conteo, y luego confirmar.
    \item Verificar el radicado, reintentar y descargar la constancia.
  \end{enumerate}
  \medskip Ensayo previsto de hasta diez minutos. Integración pendiente.
\end{slide}

\begin{slide}{Diseño en n8n}
  \begin{block}{Ruta de conversación}
    Webhook de chat $\rightarrow$ agente con temperatura 0 y memoria por sesión $\rightarrow$ consulta de pedido y política $\rightarrow$ evaluación y borrador.
  \end{block}
  \begin{block}{Ruta de radicación}
    Webhook de confirmación $\rightarrow$ verificación de sesión y borrador $\rightarrow$ escritura idempotente $\rightarrow$ radicado persistido.
  \end{block}
  \medskip Streamlit accede a n8n por su dirección interna. PostgreSQL almacena la memoria, los datos de referencia y los casos.
\end{slide}

\begin{slide}{Del SPEC al código}
  \begin{itemize}
    \item \textbf{SPEC} fija historias y criterios medibles.
    \item \textbf{PLAN} fija contratos, topología y fuente de datos.
    \item \textbf{TASKS} vincula cada tarea con una prueba y su evidencia.
  \end{itemize}
  \begin{block}{Decisiones de diseño}
    Los plazos se calculan con código. PostgreSQL permite cargar datos reproducibles e imponer unicidad sobre el borrador confirmado.
  \end{block}
\end{slide}

\begin{slide}{Grounding y controles}
  \begin{itemize}
    \item Cada evaluación identifica el pedido, la regla, la versión y las fechas utilizadas.
    \item El cliente exige una casilla. El servidor valida el borrador, la sesión y la confirmación antes de crear el caso.
    \item Un pedido desconocido o una regla insuficiente genera aclaración o revisión humana. Los errores de servicio se muestran en español.
  \end{itemize}
  \medskip Evidencia prevista en el JSON crudo, las ejecuciones de n8n y el registro de casos.
\end{slide}

\begin{slide}{Evaluación}
  \begin{itemize}
    \item PLN: corpus etiquetado, macro-F1, confusiones y abstenciones.
    \item Reglas: último día, festivo, plazo vencido y evidencia incompleta.
    \item Sistema: dos sesiones, escritura bloqueada, confirmación, reintento y constancia.
  \end{itemize}
  \begin{block}{Estado de este borrador}
    Existe una prueba de humo del entorno con el agente de clase. Las pruebas de Jaleco se incorporarán cuando estén ejecutadas.
  \end{block}
\end{slide}

\begin{slide}{Limitaciones y trabajo futuro}
  \begin{itemize}
    \item Demostración con tienda y pedidos sintéticos.
    \item Cobertura limitada a las reglas y excepciones documentadas.
    \item Validar el ahorro de tiempo con más casos e incorporar datos operativos autorizados.
    \item Explorar integraciones de soporte y transporte tras evaluar el flujo inicial.
  \end{itemize}
\end{slide}

\begin{slide}{Equipo, repositorio y demo}
  \begin{itemize}
    \item Frentes de trabajo: datos y evaluación, n8n, Streamlit, SDD y presentación.
    \item Reparto por integrante y enlace de demo pendientes de integración.
    \item \href{https://github.com/juancohen117/ProcesamientoLenguajeNatural/tree/main/Proyecto-Final-Grupo5}{Repositorio del Grupo 5 en GitHub}.
  \end{itemize}
  \medskip Presentación inicial para revisión del equipo.
\end{slide}

\end{document}
